\subsection{Resources Managed} \label{resources-managed} % done

FreeBSD manages system resources such as the underlying CPU, main memory, pipes, sockets, files, and disk devices
\parencite{mckusick2014}. System resources such as pipes, sockets, files, and disk devices can be accessed by the
\textit{descriptor array}, which is one of the kernel resources FreeBSD manages aside from credentials and signal states
\parencite{mckusick2014}. Each process is assigned to its own descriptor array listing down the devices it can
use \parencite{mckusick2014}. Each thread is assigned to its own set of registers, program counter, and stack \parencite{mckusick2014}.
CPU and memory are the most preemptable resources while accessing serial I/O ports cannot be preempted and requires
exclusive access. Comprehensive resource tracking and management in FreeBSD is handled by the \texttt{rctl} command
\parencite{wojtczak2022}. Configurations such as \texttt{cputime}, \texttt{memoryuse} for resident set size,
\texttt{vmemoryuse} for virtual address space, \texttt{swapuse} for swap usage, \texttt{openfiles} for open file
descriptors, \texttt{nthr} for the number of threads used, and \texttt{shmsize} for \textit{SysV shared memory segments}
\parencite{wojtczak2022}. These settings can be adjusted for each process, user, or jail. This granularity allows
for a fine-grained control over shareable resources such as the main memory and the descriptor array and nonshareable
resources such as CPU cores and exclusive device handles. FreeBSD treats the main memory as its most critical resource,
hence FreeBSD provides mechanisms to ensure that main memory is properly utilized. For instance, it adopts
a generic \textit{"VM object"} for unifying different forms of memory storage such as swap-backed, file-backed, or physical
device-backed storage into a unified buffer cache \parencite{freebsd_vms}. This can potentially prevent redundant
memory use across the filesystem, which can avoid disk I/O and process stalls. Another interesting detail about FreeBSD
is that it treats interrupt handling as its own resource management concern. In FreeBSD, interrupts threads run at
real-time kernel priority and have their own thread context \parencite{freebsd_genarchi}. This allows interrupt threads
to block upon encountering locks \parencite{freebsd_genarchi}. However, since multiple handlers can share the same interrupt
thread, they should not sleep to prevent starvation of other interrupt handlers \parencite{freebsd_genarchi}. Despite this,
interrupt threads being able to block on locks and have their own context can improve the latency of the system.

In NetBSD, kernel entities are able to share resources associated with a process \parencite{williams_scheduler}.
Each kernel entity has its own running state \parencite{williams_scheduler}. The kernel can preempt a thread and transfer
the CPU control to the uptime handler \parencite{williams_scheduler}. With this design, NetBSD treats CPU time as a
preemptable resource and address spaces and file descriptor tables as shareable resources accessible by multiple threads
of the same process. NetBSD organizes resource management in \texttt{/sys/kern}, where it contains mechanisms for
managing signal delivery, resource control, and process tracking \parencite{netbsd_src}. Block and character devices and
descriptor tables which are covered by the I/O subsystem are treated as first-class kernel resources by NetBSD.
Memory in NetBSD is managed as a dynamically allocated and on-demand paged resource and serves as the most critical resource
in NetBSD. This is reflected in how NetBSD's virtual memory manager, \textit{UVM}, treats memory. UVM treats the main
memory as a set of \texttt{uobj} objects, which are a sequence of contiguous virtual memory regions backed in a storage space
such as a vnode or the swap \parencite{netbsd_internals}. The 5.0 release of NetBSD further extends its resource
management system to \textit{SMP} systems with processor and CPU sets for thread affinity \parencite{netbsd_5_0}.
Almost all core kernel subsystems are redesigned to adopt highly concurrent algorithms \parencite{netbsd_5_0}.
A flaw to this redesign is that it can lead to systemwide freezes under workloads with frequent switching between high CPU
and I/O operations on security subsystems caused by locking mechanisms. Nevertheless, this implementation treats
CPU cores as a dynamically assignable resource.

FreeBSD and NetBSD both treat memory as their most critical resource. Both operating systems also converge on unifying
memory storage abstractions. FreeBSD exposes
pipes, sockets, files, and devices through a per-process descriptor array alongside credentials and signal states.
This provides a fine-grained and configurable tracking across processes, users, and jails, which is done using the
\texttt{rctl} command. NetBSD similarly treats block and character devices and descriptor tables as first-class kernel
resources and organizes them in \texttt{/sys/kern}. However, they diverge in the scope and structure involved in
managing these resources. FreeBSD uses a generic VM object to combine swap-backed, file-backed, and device-backed storage
into a single unified buffer cache, while NetBSD's UVM treats memory as a set of \texttt{uobj} objects which are backed
by a vnode or swap. A key tradeoff of the presented approaches lies in how these are reflected in concurrent workloads.
NetBSD's 5.0 release incorporates highly concurrent algorithms and support for SMP systems with processor and CPU sets
for thread affinity by reworking almost all of the core kernel subsystems. This comes at a cost of potential system-wide
freezes on workloads that frequently switch between CPU and I/O loads on security subsystems. FreeBSD's jail-aware
\texttt{rctl} model offers more granularity in administrative control and more organized resource accounting. This design
is certainly well-suited to multi-tenant and containerized environments, which is consistent of using jails as FreeBSD's
primary virtualization mechanism \parencite{freebsd_jails}.

\subsection{Resource Allocation Strategies} \label{resource-allocation-strats} % done

Dynamic CPU allocation in FreeBSD is managed by the \textbf{ULE scheduler}, which is FreeBSD's primary CPU scheduling algorithm
\parencite{mckusick2014_process}. ULE enforces both fairness and priority in scheduling processes \parencite{mckusick2014_process}.
This is done by finding a stable equilibrium from by assigning a proportional CPU time depending on the nice value instead
of disproportionally allotting CPU time to the highest priority process \parencite{robserson_ule}.
ULE achieves this with per-CPU run queues and CPU topology model \parencite{robserson_ule}.
The ULE scheduler has different scheduling mechanisms for interactive and non-interactive threads. For interactive threads,
their priority numbers are gradually boosted to improve the system's responsiveness \parencite{mckusick2014_process}. While
for non-interactive threads, they are added to a time-sharing queue \parencite{mckusick2014_process}. This allows non-interactive
threads to run at least once per pass regardless of priority, which ensures that CPU access is fairly distributed across
non-interactive threads \parencite{mckusick2014_process}. This scheduling mechanism achieves a good balance of throughput
for large batched workloads and latency for interactive workloads. However, this mechanism can potentially starve low-priority
background jobs. Allocating other system resources such as memory and file resources to users are performed with login classes,
which are defined in \texttt{/etc/login.conf} \parencite{freebsd_limiting_users}. A login class is assigned to each user,
wherein it specifies caps on CPU time, memory use, file descriptor use, and process count for each user \parencite{freebsd_limiting_users}.
They are typically set statically on login time and cannot be adjusted automatically.
However, disk quotas are not managed in the login class \parencite{freebsd_limiting_users}. Instead, they are handled separately
with the \texttt{edquota} command \parencites{freebsd_limiting_users}{freebsd_filesys_quotas}.
FreeBSD employs limits on using certain resources with \textit{resource containers} or \textit{RCTL} \parencite{freebsd_resource_limits}.
The RCTL framework enforces limits on specific resources using rules \parencite{freebsd_resource_limits}. Each rule specifies
a subject, which can either be a user, process, or jail, a resource to be used, and an action, which includes \texttt{deny}
or denying an allocation, \texttt{sighup} or sending the \texttt{SIGHUP} signal to the offending process, or \texttt{log}
or logging a message to the system log \parencite{freebsd_resource_limits}. This allows for a dynamic control during runtime allocation
which can be updated without having to reboot. This also limits cascades from jails to individual processes. However, RCTL
must be explicitly enabled by setting \texttt{kern.racct.enable} to \texttt{1} \parencite{freebsd_resource_limits}. This
can only be tuned at boot time. It is also disabled by default because of the overhead that comes with it.
In reclaiming resources, specifically memory, FreeBSD's VM system makes use of the \textbf{multi-queue LRU model} to reclaim memory
by evicting specific pages \parencite{dillon}. It evicts dirty pages when memory is stressed \parencite{dillon}. \textcite{dillon}
noted that selecting a wrong page to be evicted can cost the system hundreds of thousands of CPU cycles. Hence significant
overhead might be required to make the correct choices in selecting pages to be eviceted.

Thread allocation in NetBSD is mainly governed by two schedulers that can be selected at compile time, which are the \textit{4.4BSD
	scheduler} and the \textit{M2 scheduler} \parencite{netbsd_internals}. While both schedulers both utilize a time-sharing queue
and performs dynamic priority calculation, they differ in how the priorities are recalculated. The 4.4BSD scheduler recalculates
priority using a digital decay filter which checks the recent CPU utilization of a thread on \texttt{l\_estcpu} and adjusts
the priority in each second \parencite{netbsd_internals}. However, the M2 scheduler adopts a priority recalculation algorithm
based on SVR4 or Solaris-style time quantum logic, which can perform priority recalculation in $O(1)$ time \parencite{netbsd_internals}.
While the M2 scheduler provides more interactivity in some workloads, it is still not ready for production use. For \textit{symmetric
	multiprocessing (SMP)} systems, the scheduler dispatcher allocates CPU to these systems with per-CPU run queues and heuristic-based
balancing \parencite{rasiukevicius2009}. Thread affinity can also be set explicitly to reduce \textit{cache thrashing} by
setting explicit bindings of threads to processors \parencite{rasiukevicius2009}. In performing per-process resource allocation,
it has to go through the POSIX \texttt{setrlimit()} and \texttt{getrlimit()} interface \parencite{netbsd_setrlimit}. System resource
consumption limits can be accessed and adjusted by calling \texttt{getrlimit()} and \texttt{setrlimit()} respectively \parencite{netbsd_setrlimit}.
Limits on resource consumption are further specified by their soft and hard limits. A \textit{soft limit} is the set limit that
is currently enforced by the system while a \textit{hard limit} is the permitted ceiling limit for consuming a resource
\parencite{netbsd_setrlimit}. Adjusting hard limits is only restricted to the superuser while unprivileged users can adjust
the soft limits of a particular resource \parencite{netbsd_setrlimit}. Infinite limits can also be specified with the
\texttt{RLIM\_INFINITY} macro \parencite{netbsd_setrlimit}.

The two operating systems mainly differ in how resource limits are imposed and the maturity of their scheduling policies.
FreeBSD relies on a production-grade ULE scheduler that enforces both fairness and priority by assigning CPU time proportional
to the assigned nice value. NetBSD offers two scheduling algorithms, with a production-ready but older 4.4BSD scheduler
that performs priority recalculation based on recent CPU utilization and a faster but less mature M2 scheduler that is
able to perform priority recalculation in $O(1)$ time. FreeBSD's mature single-scheduler approach prioritizes stability
and balanced throughput-latency relationship at the risk of starving low-priority background processes, whereas NetBSD's
dual-scheduler approach offers flexibility. Moreover, in enforcing resource limits, FreeBSD combines statically assigned
login classes with dynamically updatable RCTL framework, albeit RCTL is only available for developmental environments
due to the significant overhead it brings. This is in contrast with NetBSD that provides routing to per-process limits
through the POSIX \texttt{setrlimit()}/\texttt{getrlimit()} interface with soft and hard limits. This approach favors
portability and POSIX compliance over FreeBSD's dynamic but heavier runtime enforcement.

\subsection{Resource Tracking} \label{resource-tracking} % done

In FreeBSD, resource accounting is done at the kernel level through \textbf{RACCT (Resource Accounting)}, which was first introduced
in FreeBSD 9.0 \parencite{freebsd_9_relnotes}. RACCT maintains a per-process, per-jail, and per-login class accounting data which is
used for resources such as CPU times, resident memory use, and the number of open files \parencite{freebsd_racct_status}. It operates
together with the RCTL framework previously discussed in Section \ref{resource-allocation-strats} to enforce limits and trigger actions
that can be configured whenever the enforced limits are exceeded \parencite{freebsd_racct_status}. It is required that the accounting
routines will be invoked by RACCT on every allocation and deallocation, it brings overhead to the system \parencite{freebsd_racct_status}.
Because of this, it is disabled by default and system administrators must explicitly set \texttt{kern.racct.enable} to \texttt{1} in
\texttt{/boot/loader.conf} to enable it \parencites{freebsd_racct_status}{freebsd_resource_limits}.
RACCT has two primary primitives in tracking how resources are use, which are the \texttt{racct\_add()}, \texttt{racct\_sub()},
and \texttt{racct\_set()} functions. \texttt{racct\_add()} and \texttt{racct\_sub()} update the usage counters on a resource allocation
or release respectively, while \texttt{racct\_set()} is used to manually set the usage counter to a specific value \parencite{freebsd_resource_limits}.
This function is typically used for file descriptors, wherein file descriptor tables can be shared between processes. Enforcing an incremental
tracking system would be imprecise \parencite{freebsd_resource_limits}. As a consequence, file-descriptor accounting is documented as
somewhat imprecise, but not critically imprecise \parencite{freebsd_resource_limits}.

Resource reclamation in FreeBSD is a two-phase event which is tied to terminating a process.
When a process is terminated, the \texttt{exit()} routine is called to kill all threads associated with that process \parencite{mckusick2014_term}.
Additionally, it also resets the execution state on kernel mode, removes the process from the \texttt{allproc} list, and moves that process
to the zombie list specified by \texttt{zombproc} \parencite{mckusick2014_term}. The accumulated resource usage is bundled into a
\texttt{p\_ru} structure which is retained for the parent to be retrieved by calling the \texttt{wait4()} function \parencite{mckusick2014_term}.
The kernel then reclaims all descriptors that were in use. If the terminating process holds the final reference to an object such as a file or a
socket, the object's manager is notified to perform some final cleanup such as deleting the file or deallocating the socket
\parencite{freebsd_design4bsd}. For performance-sensitive workloads wherein always-on accounting is too expensive, FreeBSD provides a profiling
tool called \textbf{DTrace} that can be used by the kernel and userland to expose resource usage at production granularity when enabled
\parencite{lavigne2021}.

In NetBSD, it tracks per-process resource usage with \texttt{proc.pid.rlimit} \texttt{sysctl} MIB hierarchy, wherein the resource limits
of each process are exposed, which can include CPU time, virtual memory, file size, and descriptor count such as queryable and settable
\texttt{sysctl} nodes \parencite{netbsd_sysctl}. The \texttt{proc.curproc} shorthand targets the current process, while arbitrary processes
can be addressed by their PID \parencite{netbsd_sysctl}. Instead of tracking usage counters adopted in FreeBSD, NetBSD tracks the configured
limits instead. This implementation makes it lighter but can be less capable for performing aggregate accounting tasks across multiple jails
or user classes. The limits for each process are read and modified through the POSIX \texttt{getrlimit()} and \texttt{setrlimit()} interface
mentioned in Section \ref{resource-allocation-strats}. This interface is integrated with separate soft and hard limits \parencite{netbsd_getrlimit}.
\textcite{netbsd_getrlimit} points out a bug wherein physical memory size limits can be bypassed when the resident set size limit
\texttt{RLIMIT\_RSS} is not implemennted in UVM. Resource usage is charged at termination by the \textbf{Common Scheduler Framework (CSF)},
wherein a process's accumulated usage can be charged to its parent on exit with the \texttt{sched\_proc\_exit()} hook \parencite{netbsd_csf}.
In addition, a periodic \texttt{sched\_pstats()} routine updates the usage statistics of a process and verifies its CPU usage \parencite{netbsd_csf}.
For memory usage, UVM reclaims pages by flushing the dirty pages through their backing stores by performing the \textit{putpages} operation,
when closing files, or when explicit filesystem synchronization occurs \parencite{netbsd_internals}

The two operating systems adopt opposing philosophies in terms of resource tracking. FreeBSD adopts an \textit{active accounting} model
wherein live usage counters are regularly tracked and maintained. This enables hierarchical aggregation across processes, users, and jails.
In addition, this also enforces resource limits at runtime. However, this comes at a cost of overhead for each allocation, justified by its
disabled-by-default posture. In contrast, NetBSD incorporates a much lighter approach with its \textit{limit-oriented} model which records
configured limits instead of live usage. This avoids overhead that comes with accounting, but it has no built-in aggregation across jails or
login classes. The presence of DTrace provides FreeBSD an additional opt-in path for observability with no idle cost, whereas NetBSD's
approach is simpler and requires minimal overhead by exposing limits through \texttt{sysctl}. FreeBSD's RACCT requires incurring significant
costs that comes with richer enforcement of resource usage and aggregation, while NetBSD remains lean but cannot provide fine-grained management
for resource attribution.

\subsection{Load Balancing and Distribution} \label{load-balancing-distribution} % done

The ULE scheduler previously mentioned in Section \ref{resource-allocation-strats} primarily governs CPU load-balancing in FreeBSD by coupling per-CPU
run queues with two complementary mitigation strategies. When a CPU becomes idle, FreeBSD performs a \textit{pull migration} wherein
\texttt{tdq\_idled()} looks for other CPUs whose work can be stolen \parencite{mckusick2014_process}. A long-term load balancer also performs
\textit{push migration} by moving a thread from the busiest CPU to the least busy CPU once per second
\parencites{mckusick2014_process}{mckusick2020}. For a CPU's work to be stolen, a load threshold must be exceeded to
prevent excessive migration and cache thrashing \parencite{mckusick2020}. Thread affinity and processor utilization
are treated by ULE as differing goals that need to be addressed during placement. \texttt{sched\_pickcup()} selects the best CPU when a
thread can be executed by balancing thread affinity, current load, and topology relationships such as cache sharing and NUMA proximity
\parencite{robserson_ule}. This is evaluated on three events: a thread becomes runnable, a CPU goes idle, and the periodic rebalancer
ticks \parencite{mckusick2014_process}. This avoids overhead that comes with constant recalculations, however, it can leave short-term
imbalances between ticks \parencite{mckusick2014_process}. Furthermore, ULE treats symmetric multithreading as a special case of NUMA,
wherein scheduling two threads from the same process onto related logical CPUs are prevented unless necessary with CPU topology information
provided by machine-dependent code \parencite{roberson_ule}. This can improve cache efficiency and reduce contention between threads sharing
the same physical execution units \parencite{robserson_ule}. When compared against Linux's CFS, it was found by \textcite{bouron2018}
that ULE achieves a better long-term load balancing when compared to CFS. However, ULE might cause temporary starvation on identical threads,
which can apparently improve throughput on some workloads by concentrating CPU time on active threads \parencite{bouron2018}.

In NetBSD, load balancing and distribution especially on SMP systems is done using a dispatcher that also utilizes per-CPU run queues with
thread migration and heuristic-based balancing \parencite{rasiukevicius2009}. This attempts to reconcile the same opposing goals of
increasing processor utilization and maintaining thread affinity. The dispatcher's decisions are described as potentially suboptimal
because of its heuristic nature instead of adopting a more deterministic and topology-aware algorithm \parencite{rasiukevicius2009}.
NetBSD 5.0 introduced \textit{processor sets} and \textit{dynamic CPU sets} to provide administrators with explicit control over how threads
and processes are bound to specific CPU groups \parencite{netbsd_5_0}. A processor set follows a resource-pool model wherein bound
processors can only run on their assigned threads, which allows critical workloads to reserve CPU resources with an added risk of
putting those CPU resources on idle when the bound threads gets blocked \parencite{netbsd_5_0}. In terms of memory, NetBSD's UVM
pagedaemon distributes reclamation pressure across page types such as anonymous, executable, and file-cache pages using six configurable
\texttt{sysctl} parameters that define a lower and upper limit per page type \parencite{selonen}. When a page type falls below its
lower limit, the pagedaemon reactivates the pages instead of freeing them \parencite{selonen}.

FreeBSD and NetBSD both use per-CPU run queues and frame distribution in balancing affinity and utilization. However, they largely differ
in their determinism and granularity. FreeBSD's ULE draws on machine-dependent topology data and combines event-driven pull and push
migrations, making ULE topologically-aware. This approach is able to balance the two goals in the long term. However, this comes at a
cost of brief imbalances between ticks and possible starvation. While NetBSD's dispatcher approach is not exactly optimal, this
disadvantage is compensated by exposing explicit processor and CPU sets which provides more deterministic control over placement.
This contrast also seems to appear in how memory is distributed. FreeBSD relies on the multi-queue LRU reclamation of its VM system
as described in Section \ref{resource-allocation-strats}, whereas NetBSD's per-type pagedaemon balancing is more fine-grained but
difficult to tune.

\subsection{Handling Deadlocks} \label{handling-deadlock} % done

Both FreeBSD and NetBSD primarily incorporate \textit{deadlock prevention} in dealing with deadlocks. This is done by enforcing a
consistent global lock ordering instead of detecting at runtime. They mainly differ in the maturity of the tooling for performing
such mechanism. This discipline is enforced by the idea that if every thread acquires locks in the same order, the circular-wait
condition cannot arise \parencites{baldwin2002}{netbsd_mutex}.

In FreeBSD, this is done with a dynamic lock-order verifier called the \textbf{witness} \parencite{baldwin2002}. The witness builds
a directed graph of lock-acquisition relationships at runtime \parencite{baldwin2002}. For instance, if a thread acquires lock B
while holding lock A, the witness records this relationship. If another thread later tries to acquire the same pair in reverse
order, the witness logs a warning and optionally drops into the kernel debugger \parencite{baldwin2002}. The \texttt{WITNESS}
kernel option maintains this graph by lock type and checks for potential cycles at each acquisition \parencite{freebsd_kerneldebug}.
When a cycle is detected, the system generates a warning and stack trace to the console \parencite{freebsd_kerneldebug}. Additionally,
the \texttt{WITNESS\_SKIPSPIN} option skips the frequently acquired spin locks, which can potentially reduce overhead
\parencite{freebsd_kerneldebug}. Because of the extremely high performance cost it comes with, witness is disabled in production
kernels and is only accessible in development environments \parencite{baldwin2002}. FreeBSD also prevents specific deadlock classes
structurally. For instance, a spin mutex can disable interrupts on the local CPU while being held so that another interrupt cannot
attempt to reacquire the same mutex being held by the interrupted thread \parencites{mckusick2014_sync}{freebsd_mutex}. This, however,
comes at the cost of increased interrupt latency \parencites{mckusick2014_sync}{freebsd_mutex}. FreeBSD also incorporates an
additional role that distinguishes \textit{bonded sleeps}, which blocks acquisitions for a block when it is currently owned by
a specific process, from \textit{unbounded sleeps}, which waits for an external event \parencite{freebsd_locking}. A thread in a
bounded sleep must never wait on a lock being held by a thread in an unbounded sleep \parencite{freebsd_locking}. Moreover,
recursive \texttt{Giant} mutex protecting data structures are dropped during unbounded sleeps to prevent cascading blocks
\parencite{freebsd_locking}.

NetBSD adopts the same discipline for preventing deadlocks and adds onto it by enforcing a strict hierarchy wherein adaptive mutexes
and sleep locks cannot be acquired whenver a spin mutex is being held \parencite{netbsd_mutex}. This approach can prevent a class
of interrupt-handler deadlocks \parencite{netbsd_mutex}. The \texttt{LOCKDEBUG} compile-time option can be set during development
to record lock-state information for diagnosing deadlocks \parencite{netbsd_lock}. \texttt{LOCKDEBUG} is disabled for production
environments due to the heavy overhead it comes with. For simplelocks, it converts a thread attempting
to relock a lock it already holds into a \textit{locking against myself} similar to FreeBSD's witness \parencite{netbsd_lockmgr}.
NetBSD plans on implementing a deadlock detection mechanism, but it still remains as an open project
\parencite{netbsd_deadlock_detector}.

The practical consequences of relying on prevention is evident in NetBSD's recent jails development, wherein a full system freeze
can be caused from \texttt{secmodel\_eval} calls on UVM and scheduling paths under mixed high CPU and I/O loads
\parencite{petermann_jails}. This is rooted to lock-ordering violations, recursion, and subtle lock interactions at security-model
integration points \parencite{petermann_jails}. As of now, NetBSD does not currently adopt a recovery mechanism, hence reboots
are often done to resolve this. But overall, both systems share a prevention-first philosophy for handling deadlocks with
comparable compile-time debugging tollings. However, FreeBSD's runtime witness verifier provides stronger in-production safeguards.
In contrast, NetBSD's prevention-only model leaves complex codepaths exposed to unrecoverable freezes documented by
\textcite{petermann_jails}.

\subsection{Resource Contention and Resolution} \label{resource-contention} % done

When multiple threads compete for a short-term lock such as a mutex, reader-writer lock, or read-mostly lock, FreeBSD tracks both
the lock owner and waiting threads in a \textbf{tunstile} data structure \parencite{mckusick2014_sync}. The waiting threads
are sorted within the turnstile by priority. FreeBSD also enforces a form of \textit{priority inheritance} wherein priority is
recursively propagated to the owner if a waiter has a higher priority than the current lock holder \parencite{mckusick2014_sync}.
The tunstile state itself is allocated from a pool when a contention do occur since a thread can sleep on one mutex at a time
\parencite{mckusick2014_sync}. This ensures that the number of simultaneously active turnstiles never exceeds the number of
threads \parencite{mckusick2014_sync}. This also controls potential memory overhead fluctuations \parencite{mckusick2014_sync}.

At the scheduling level, the ULE scheduler prevents CPU starvation among threads with the same priority with time slices
\parencite{mckusick2014_process}. A time slice sets the maximum latency of a thread of a given priority before running
\parencite{mckusick2014_process}. Then, the ULE dynamically adjusts the time slice depending on the current system load while
ensuring that it does not go beyond a minimum value to avoid thrashing \parencite{mckusick2014_process}. A round-robin approach
is then adopted for cycling through these threads with the same priority under high contention \parencite{mckusick2014_process}.
This approach is in contrast to the traditional BSD scheduler, which makes use of a single global run queue and a single global
scheduler lock, which becomes a contention bottleneck as the number of runnable threads increases \parencite{mckusick2014_process}.
This often happens since priority calculations need to iterate over every runnable thread while holding highly contended locks,
blocking forks, exits, and context switches at the same time in the process \parencite{mckusick2014_rpcoess}. For lock-level
contention, FreeBSd leans on \textit{lent priority} in bounded sleeps \parencite{freebsd_locking}. In this approach, a thread
currently waiting on a mutex lends its priority to the owner to allow the owner to be prioritized over low-priority threads,
which allows the lock to be released sooner \parencite{freebsd_locking}. For unbounded sleeps, blocked threads are awakened
in FIFO, which enforces fairness but not priority \parencite{freebsd_locking}. This leads to potential starvations on low-priority
threads under very high contentions.

In handling lock contentions, NetBSD makes use of \textbf{adaptive mutexes}, which behaves similarly to spinlocks that sleeps
whenever the expected hold time is long \parencite{netbsd_locking}. This busy-waiting approach is much preferred considering
that the system encounters short hold times more frequently \parencite{netbsd_locking}. Unlike a traditional spinlock, a thread
holding an adaptive mutex can be preempted by the kernel, which can hurt latency for real-time applications. At the scheduler level,
NetBSD's time-sharing scheduler prioritizes I/O-bound threads upon waking up from an I/O operation, which preserves fairness and
increases interactive responsiveness \parencite{rasiukevicius2009}. Each \textit{light-weight process (LWP)} carries a dynamically
calculated priority value and time quantum \parencite{rasiukevicius2009}. \texttt{SCHED\_FIFO} threads have fixed priorities
and are executed till completion, while \texttt{SCHED\_RR} threads are scheduled in a round-robin fashion within their class to
prevent starvation \parencite{rasiukevicius2009}.

Both designs share lazy contention-tracking allowing lightweight locks to resemble turnstiles under heavy contention. Both
also incorporates priority inheritance and round-robin cycling to limit starvation across threads of the same priority. However,
they differ in behavior under mixed CPU and I/O loads. While NetBSD's contention resolution is sufficient on normal workloads,
it can result to full system freezes when locking interactions in security-adjacent subsystems such as \texttt{secmodel\_eval}
are performed \parencite{petermann_jails}. This highlights an insufficient isolation between hot subsystem paths and the security
enforcement layer. FreeBSd's more mature approach of fine-grained locking per subsystem provides a much stronger isolation,
albeit the complexity and overhead it brings.

\subsection{Performance and Efficiency} \label{resource-performance} % done

Enforcing efficient resource management in FreeBSD requires the replacement of a single \texttt{Giant} lock with fine-grained
locking mechanism through the \textit{SMPng} project, which was completed in FreeBSD 6.0 \parencite{freebsd_smp}.
The \texttt{Giant} lock prevented concurrent execution of kernel threads, which causes high contentions and minimal SMP speedups
on kernel-intensive workloads such as networking and filesystem I/O \parencite{freebsd_smp}. A more fine-grained locking
mechanism prevents such limitations, but brings its own set of tradeoffs by increasing per-packet and per-syscall processing
cost through synchronization, which requires active tuning \parencite{freebsd_smp}. At the scheduling layer, the ULE scheduler's
event-driven design achieves a good balance of efficiency and fairness \parencite{roberson_ule}. This balance is achieved by
penalizing threads that consume their entire time slice and boosting waiting threads that do not consume their entire time slice
\parencite{robserson_ule}. This balance is also achieved by assigning smaller time slices to high-priority and interactive
threads and assigning large time slices to low-priority and CPU-bound threads \parencite{roberson_ule}. This efficiency has
has been studied by \textcite{bouron2018} by conducting benchmarks on the ULE. \textcite{bouron2018} found that
\texttt{sched\_pickcpu()} may scan all CPU cores by up to three times when selecting a placement, which results in consuming
$13\%$ of the total CPU cycles on scanning alone. It has also been found that the ULE has a tendency of starving background
threads when threads classified as background threads receive no CPU time in the benchmarks conducted by \parencite{bouron2018}.
While this boots the interactive performance of the system, it is at a cost of poor enforcement of fairness.

Both FreeBSD and NetBSD adopted the same approach in their user-space memory allocator. They replaced the legacy \texttt{phkmalloc}
allocator with a newer \textbf{jemalloc} allocator because \texttt{phkmalloc} is unable to scale in multithreaded workloads
\parencite{evans2006}. It is found that increasing the number of threads corresponds with a decrease in overall allocation
throughput with \texttt{phkmalloc} \parencite{evans2006}. Whereas jemalloc is able to increase allocation throughput with more
threads with its per-arena design \parencite{evans2006}. FreeBSD incorporated jemalloc in the 7.0 release, while NetBSD adopted
jemalloc as its default user-space allocator in the 5.0 release \parencites{netbsd_jemalloc}{evans2006}. Scaling improvements
were added in FreeBSD 7.0 by replacing shared locks, which is a post-SMPng bottleneck notably on multithreaded MySQL workloads,
with a sleepable mutex primitive \parencite{kennaway}. As for NetBSD, efforts to significantly improve its efficiency on the
5.0 release are reflected in overhauling a majority of the kernel subsystems, including virtual memory, memory allocations,
and filesystem frameworks \parencites{lwn330994}{lwn332022}. This improvement replaced shared coarse-grained locking to
fine-grained concurrent algorithms and allowing multiple processors to run kernel code simultaneously \parencites{lwn330994}{lwn332022}.
NetBSD's adaptive mutexes embody a similar balance between efficiency and correctness, which defaults to busy-waiting when hold times
are expected to be short to avoid context-switching overhead in the common case \parencite{netbsd_locking}. However, this can
result in potentially spending a lot of CPU time spinning under high contention before going back to sleeping \parencite{netbsd_locing}.
A case study conducted by \textcite{brown_seltzer} explored NetBSD's performance across Intel x86 generations and found that
off-chip memory-system design significantly influences the performance of the operating system. Furthermore, it has been found
that some kernel design decisions for dealing with cache and memory-bus constraints do not always transfer between chip
generations \parencite{brown_seltzer}.

Both operating systems share the same techniques of more fine-grained and per-subsystem locking and scalable concurrent
allocator. However, these approaches mainly differ in when they are implemented and the emphasis of these design decisions.
FreeBSD adopted this approach earlier with its transition to SMP systems with SMPng which is then further pushed into individual
subsystems. While this significantly aid in boosting throughput on kernel-intensive workloads, this approach incur a lot of
synchronization cost per operation. NetBSD's 5.0 overhaul also achieved comparable fine-grained concurrency and real-time
responsiveness. However, NetBSD is still limited on its philosophy on high portability, which adds a cross-architecture tuning
burden. Moreover, this contention previously discussed in Section \ref{resource-contention} leaves an edge case causing
system freezes on high workloads.

\subsection{Scalability} \label{scalability} % done

FreeBSD's scalability on SMP systems has been refined over its releases. It was initially introduced in the 5.x releases
which began replacing the \texttt{Giant} lock \parencite{freebsd_vs_dragonfly}. In FreeBSD 7.x releases, the ULE scheduler
was introduced along with significant lock-granularity improvements \parencite{freebsd_vs_dragonfly}. This is further extended
in FreeBSD 10 by adding per-structure locks, reader-writer locks, and lockless algorithms to most subsystems
\parencite{freebsd_vs_dragonfly}. A tipping point has reached in FreeBSD 14 with the addition of NUMA awareness
\parencite{freebsd_vs_dragonfly}. The documented cumulative results shows how FreeBSD scaled to more than 100 cores for
various workloads \parencite{freebsd_vs_dragonfly}. It is important to note that this documented result is heavily
workload-dependent. A study conducted by \textcite{cui2011} comparing Linux and FreeBSD in AMD 32-core systems found that
FreeBSD scales better in workloads that perform intensive socket operations while Linux scales better for file-descriptor
and process-intensive operations. \textcite{cui2011} also found that the synchronization primitives which protect the shared
kernel data structures serve as the primary bottleneck in inhibiting the scalability of the system \parencite{cui2011}.
Specifically, \textcite{freebsd_processor_trends} identified a NUMA-related ceiling in network I/O, wherein a dual-socket
system achieved a consistent $35 G/s$ throughput only when the sender, receiver, and the texttt{netisr} thread were all scheduled
on the same NUMA socket. When moving to a dual-socket configuration, this design introduced significant latency degrading
the usable bandwidth \parencite{freebsd_processor_trends}. This issue seems to be caused by using a single \texttt{netisr}
thread or the entire system and performing scheduling decisions that does not consider the locality of inter-thread communications
\parencite{freebsd_processor_trends}. On the administrator side, adding or removing rules requires the kernel to walk through
all \texttt{uidinfo}, \texttt{loginclass}, jail, and process structures in the \texttt{rctl} hierarchical resource-limit
system \parencite{freebsd_resource_limits}. This means that administrative operations do not scale well with increasing users
and jails despite the fact that no global tree is managed in performance-critical paths \parencite{freebsd_resource_limits}.

NetBSD's approach to scalability began in the 5.0 release with the introduction of processor and CPU sets, and dynamic CPU
affinity as its primary resource scaling mechanisms across cores \parencite{netbsd_5_0}. The older scheduler-activation
model mentioned in Section \ref{resources-managed}, which scaled poorly on multiprocessor and real-time systems,
is replaced by a 1:1 threading model \parencite{williams_scheduler}. As discussed in Section \ref{load-balancing}, the
dispatcher balances utilization and affinity by performing heuristic decisions based on observed thread behavior
\parencite{rasiukevicius2009}. This is recognized as a known limitation to NetBSD, which is relative to ULE's explicit and
machine-dependent topology awareness \parencite{rasiukevicius2009}. Despite this known limitation, NetBSD employs other
strategies to improve the scalability of the system. Notably, NetBSD's approach to scalability focuses on expanding to a
wide range of platforms and architecture instead of scaling vertically. NetBSD is well supported on a wide range
architectures and instruction sets, which includes VAX minicomputers, modern ARM architectures, Raspberry Pi chipsets,
and RISC-V architectures \parencite{netbsd_portability}. All of these releases are built on a single source tree. This
portability-first design makes NetBSD's resource-management abstractions directly align to system architectures
\parencite{netbsd_portability}. This allows NetBSD to scale well horizontally but incorporating architecture-specific
optimizations might incur additional overhead. NetBSD is also known for its lightweight scaling, with conducted benchmarks
by \textcite{gmcgarry} showing that NetBSD offers faster socket-allocation times and a lower increase rate of allocation
times as the number of allocated sockets grows. This suggests that a lighter-weight design can offer competitive scaling
especially in network-heavy workloads.

This highlights fundamentally different approaches to scaling. On one hand, FreeBSD pursues \textit{vertical scalability}
by enabling a single system image to scale to high core counts with NUMA-aware scheduling and per-subsystem locking.
This design allows FreeBSD to excel well in networking workloads, however, faces limitations on processor-heavy workloads.
FreeBSD's vertical scaling approach also incurs costs in administrative management with \texttt{rctl}. On the other hand,
NetBSD pursues \textit{horizontal scalability} by architecturing a single codebase to support a broad range of hardware.
This approach provides NetBSD with a competitive lightweight performance especially on network primitives, despite its
performance being counteracted with its heuristic scheduler and the overhead that comes with architecture-specific
optimizations. This shows that FreeBSD favors optimization depth on hardware that supports it while NetBSD favors
breadth of portability across different platforms.

\subsection{Summary} \label{resource-summary}

Considering their similar BSD lineage, FreeBSD and NetBSD converge on a lot of resource-management techniques. For
instance, both treat memory as their most critical resource. Both unify memory storage abstractions in the form of VM objects
in FreeBSD and UVM \texttt{uobj} objects in NetBSD. Both also use per-CPU run queues with priority inheritance and
round-robin cycling to limit starvation. Both also adopted a more scalable jemalloc allocator. And lastly, both
adopt enforcing global lock ordering in deadlock prevention. They differ in the focus of their design choices. FreeBSD
focuses on the depth of their resource management. This is reflected in their adoption of a single ULE scheduler
with NUMA-aware and topologically-driven placement, RACCTL/RCTL framework for hierarchical and per-jail resource
control, a runtime witness deadlock verifier, and fine-grained per-subsystem locking through the SMPng project.
This depth also brought significant disadvantages to FreeBSD. RACCT's approach on accounting on every allocation is expensive,
hence it is disabled by default. In addition, RCTL requires explicit boot-time enabling. Witnesses are also too
expensive to be used in production kernels. Managing \texttt{rctl} also scales poorly as the number of users and jails
grow. NetBSD focuses on the breadth and leanness of their resource management. This is reflected in their light
limit-oriented tracking model exposed through \texttt{sysctl} and the POSIX \texttt{setrlimit()}/\texttt{getrlimit()}
interface, explicit affinity control through their processor and dynamic CPU sets, and support for a wide range
of architectures from VAX to RISC-V on a single source tree. However, NetBSD suffers from the significant tradeoffs
of its two primary schedulers, lack of aggregation across jails or login classes, and the possibility of unrecoverable
freezes from its deadlock prevention approach.

The tradeoffs between FreeBSD and NetBSD in resource management is explicit across the four comparison dimensions.
In terms of efficiency, while both are able to achieve fine-grained concurrency comparable to each other's implementations,
FreeBSD's ULE scheduler is able to balance throughput and latency through event-driven time-slice tuning while NetBSD's
adaptive mutexes default to busy-waiting on short holds. However, FreeBSD's \texttt{sched\_pickcpu()} can consume
a significant number of CPU cycles when performing placement scanning. NetBSD can also waste CPU spinning under high
contention. In terms of fairness, both adopt priority mechanisms for enforcing it. However, FreeBSD has the tendency
to starve background processes in exchange for high responsiveness, which has been verified in benchmarks. In terms
of scalability, both have opposing directions towards scaling. FreeBSD achieves vertical scaling with its refined
SMP support and excellent scalings on networking performance. This design, however, comes with NUMA ceilings and
poor \texttt{rctl} administrative scaling. NetBSD achieves horizontal scaling with its broad hardware support and
a competitive lightweight performance on network primitives. This design, however, comes with non-optimal performance
from its heuristic scheduling and overhead from architecture-specific optimizations. In terms of complexity,
FreeBSD clearly adopted more complex design choices and implementations with its richer enforcement, tooling,
and configuration surface. This is unlike NetBSD's simpler and uniform design choices, which lacks fine-grained
resource attribution.

The presented tradeoffs show which specific workloads and environments the two excel at. FreeBSD's hierarchical
RACCTL/RCTL resource accounting and jail-aware limits, its turnstile-based priority inheritance with dynamically-
adjusted time slices, and its ULE load balancing managed through push and pull migrations makes FreeBSD excel
in high-intensity workloads with a large number of users and processes competing for limited resources under
heavy loads. This comes at a cost of background job starvation and per-operation synchronization overhead.
Under this same heavy load, NetBSD is more likely to encounter full system freezes because of the high likelihood
of locking interactions and mixed CPU and I/O loads. In contrast, NetBSD's limit-oriented tracking, its lightweight
socket-allocation scaling, and its portability-first design makes NetBSD excel in resource-constrained
workloads with limited CPU and memory resources. FreeBSD's demanding design choices would not excel under
such environments.
